% Folder 142, batch 07 — archivists' page 121 (a single page; the folder has 121 pages).
% Pass: Opus 5.5 (claude-opus-5-5), 2026-10-02 — first pass, unchecked against the pages by a human.
% Revised: Fable 5.2 (claude-fable-5-2), 2026-11-14 — read again on the facsimile, not checked by a reader.
% Both « lacet » confirmed on an enlarged crop (#145).
%
% The leaf carries two lines at its head, under the author's page number 7;
% the rest of the sheet is blank.
\documentclass[11pt,a4paper]{article}
\input{../preamble/grothendieck}

\folder{142}
\batch{7}
\pages{121}{121}
\foldertitle{Action de Galois sur Teichmüller : notes manuscrites (s.d.).}
\dating{[à partir de 1978]}
\watermark{Édition de démonstration}

\begin{document}

\page{121}
\note{p. 7 de l'auteur. La phrase commence au milieu d'un argument qui court depuis la page précédente, hors de ce lot ; le reste du feuillet est blanc.}

% « lacet », twice: the same word both times (l, a, c, a final t), and the
% sentence (a loop around a point $R_s$, then « $\simeq 1$ ») leaves no other.
« lacet » autour d'un tel $R_s$, puis exprimer que ce
lacet est $\simeq 1$.

\end{document}
